\section{C++、HTTP 与 JSON 契约}\label{sec:res-http}

\paragraph{暴雨与雷击场景扩展}
场景生成路由接受 \fld{resilience.hazard\_type} 为 \texttt{rainstorm} 或
\texttt{lightning}，同名对象提供参数；弹性版能力接口通过
\fld{scenario\_hazards} 返回默认值、单位及范围。结果包含
\fld{hazard\_evidence} 的强度过程、支路故障概率、故障类型、等待与修复时间及模型限制。
对应 Web 恢复请求设置 \fld{respect\_fault\_windows=true}，RA 模型逐时遵守给定停运窗口，
不使用原贪心修复规则提前释放支路。此模式的维修计划由场景给定，不代表抢修队优化。
公式、字段与数值证据见本模块 \srcpath{weather\_scenarios.md}。

\subsection{C++ 入口}
\begin{center}\small
\begin{tabular}{@{}L{5.0cm}L{8.2cm}@{}}
\toprule
入口 & 语义\\\midrule
\srcpath{run\_distribution\_resilience\_assessment} & 按 \fld{model} 分派启发式、严格 MIP 或 RA 分阶段入口\\
\srcpath{run\_distribution\_resilience\_mip\_assessment} & 直接求严格多时段 hybrid AC/DC 恢复 MIP\\
\srcpath{run\_distribution\_resilience\_stage\_milp\_assessment} & 直接求 RA 分阶段拓扑与 MESS 子问题\\
\srcpath{certify\_distribution\_resilience\_dynamics} & 把已有恢复序列映射为 DAE 认证结果\\
\srcpath{run\_certified\_distribution\_resilience\_mip} & 严格 MIP 后执行一次动态二次认证；不自动反馈重求解\\
\bottomrule
\end{tabular}
\end{center}

\subsection{生产 HTTP 路由}
生产路由为 \texttt{POST /api/session/run\_distribution\_resilience}，实现位于
\srcpath{tests/run\_gui\_server.cpp}。请求前必须已有 session system；并发分析返回 HTTP 409。
典型请求为
\begin{verbatim}
{
  "model": "MultiPeriodMIPLinDistFlow",
  "mip_solver": "HiGHS",
  "horizon_hours": 24,
  "time_step_hr": 1.0,
  "mip_gap": 0.0,
  "mip_time_limit_s": 120,
  "mip_num_threads": 1,
  "allow_reconfiguration": true,
  "allow_mess_dispatch": true,
  "apply_demo_data": false,
  "faults": [
    {"domain":"DC", "branch_index":101,
     "outage_start_hr":0.0, "repair_duration_hr":6.0}
  ]
}
\end{verbatim}
\fld{apply\_demo\_data=true} 会改变请求本地系统副本，只适合演示；正式研究必须设为 false 并提交完整
网络数据。

\subsection{输入归一化}
HTTP 层接受模型与求解器别名，裁剪时域和数值范围，并解析聚合、按位置、按稳定 ID、按母线的负荷
曲线。发生多种映射时，组件级/母线级优先级必须按生产代码解释；文档示例不另造一套解析规则。
所有 profile 数值过滤 NaN/Inf 并截断为非负值。

\subsection{结果最小检查集}
客户端不得只检查 HTTP 200。至少读取：
\begin{equation}
 \{\fld{feasible},\fld{status},\fld{effective\_model},
 \fld{solver\_status},\fld{mip\_gap},\fld{model\_scope},
 \fld{validity},\fld{steps}\}.
\end{equation}
每步供电向量必须联合解释 kind/index/position/demand/served/shed；单位为 MW，聚合量乘
\fld{time\_step\_hr} 后为 MWh。优先级加权目标的单位是“优先级罚权$\times$MWh”，不是物理能量或货币。

\subsection{错误与回退}
无系统、JSON 解析失败和非法参数返回 HTTP 400；session 忙返回 409。后端不可用、time limit、MIP
无有效 incumbent 或 RA MESS 回退必须反映在状态和模型字段中。调用方不得用空 solver status 推断
最优，也不得把启发式结果标为 MIP。

\subsection{Web 指标计算扩展}
目录与独立评价接口提供书中第三章 18 项及工程运行 24 项，定义版本为
\texttt{book\_ch3\_2026.2}。默认灾害结束时刻为最后一次故障发生时刻；
显式输入优先，修复完成时刻独立保留。运行电量按完整时段积分，最后时段不丢弃；
概率指标缺少多场景后果时返回空值。公式、字段映射及非等间隔数值证据见
\srcpath{docs/modules/resilience/metrics.md}，实现为
\srcpath{src/resilience/resilience\_metrics.cpp}。
